\section{Lefschetz strings in the Chow group}\label{sec:lefschetz}

\subsection{Exact Chow blocks}\label{sec:chow-blocks}

On a pointed curve put
\[
 \pi_0=o\times C,\qquad \pi_2=C\times o,\qquad
 \pi_1=\Delta_C-\pi_0-\pi_2.
\]
They are orthogonal projectors, with
\({}^t\pi_0=\pi_2\) and \({}^t\pi_1=\pi_1\). For a partition
\([n]=S\sqcup T\sqcup U\), define
\[
 P_{S,T,U}=\bigotimes_{i\in S}\pi_{1,i}
            \otimes\bigotimes_{i\in T}\pi_{2,i}
            \otimes\bigotimes_{i\in U}\pi_{0,i}.
\]
Let \(j_T\) insert the point \([o_i]\) on factors in \(T\) and
the fundamental class \([C_i]\) on factors in \(U\). Its inverse
on this projector image is the fixed correspondence
\[
 b_T(z)=\Pi_{S*}p_{S*}
           \left(z\cdot\prod_{i\in U}H_i\right).
\]
Indeed \(b_Tj_T=\Pi_S\) and \(j_Tb_T=P_{S,T,U}\). Therefore
\begin{equation}\label{eq:chow-decomposition}
 V^r(X_E)=\bigoplus_{S,T,U}j_TU_S^{r-|T|}(E).
\end{equation}
These identities are also valid on integral images. All intersections
and pushforwards in \(b_T\) have integral coefficients.

Put \(I=[n]\setminus S\), \(m=|I|\), and
\[
 B_I=\Q[x_i:i\in I]/(x_i^2),\qquad \deg x_i=1.
\]
For \(T\subset I\), write \(a_T=\prod_{i\in T}a_i\). The
identification
\begin{equation}\label{eq:weighted-chow-block}
 u\otimes x_T\longmapsto a_Tj_Tu
\end{equation}
groups \MathEqref{eq:chow-decomposition} as
\[
 V^*(X_E)=\bigoplus_{S,c} U_S^c(E)\otimes B_I[c],
\]
where a vector in \(U_S^c\otimes B_I^t\) has codimension \(c+t\).
Multiplication by \(\ell=\sum_i a_iH_i\) becomes multiplication
by \(E_I=\sum_{i\in I}x_i\) on the second factor. On the other
factors it vanishes because \(H_i\pi_{1,i}=0\) as a Chow
correspondence. This last identity follows by composing
\(\mu_{[o_i]}=[o_i]\times[o_i]\) with the three curve projectors.

If \(z\) is integral, its coefficient \(u_T\) in
\MathEqref{eq:weighted-chow-block} satisfies
\(a_Tu_T\in\Lambda_{U_S^c}(E)\). Thus the single integer
\begin{equation}\label{eq:coordinate-denominator}
 N_a=\prod_{i=1}^n a_i
\end{equation}
clears every coefficient denominator, in every degree and over every
finite extension.

\subsection{A square-zero Lefschetz calculation}\label{sec:square-zero}

On \(B_I\), use the monomial basis \(x_T\) and define
\[
 E x_T=\sum_{i\notin T}x_{T\cup\{i\}},\qquad
 F x_T=\sum_{i\in T}x_{T\setminus\{i\}},\qquad
 Hx_T=(m-2|T|)x_T.
\]
Then \([F,E]=H\), \([H,F]=2F\), and \([H,E]=-2E\).
The operators \(E,F\) are adjoint for the positive inner product
in which the monomials are orthonormal. Let
\(*x_T=x_{I\setminus T}\).

\begin{lemma}\label[lemma]{lem:square-zero}
For \(0\le t\le m/2\),
\[
 E^{m-2t}:B_I^t\xrightarrow{\sim}B_I^{m-t},\qquad
 \ker(E^{m-2t+1}|_{B_I^t})=\ker(F|_{B_I^t}).
\]
For every rational vector space \(M\) and every
\(z\in M\otimes\ker(F|_{B_I^t})\),
\begin{equation}\label{eq:star-identity}
 E^{m-2t}z=(-1)^t(m-2t)!\,*z.
\end{equation}
\end{lemma}

\begin{proof}
The displayed commutators give the standard finite-dimensional
\(\mathfrak{sl}_2\) representation on \(B_I\). More explicitly,
if \(Fv=0\), \(v\in B_I^t\), then
\[
 FE^jv=j(m-2t-j+1)E^{j-1}v.
\]
Adjointness shows that the string generated by a nonzero such \(v\)
has nonzero terms through \(E^{m-2t}v\) and then terminates.
Decomposing successively into \(\ker F\) and \(\im E\), which
are orthogonal, gives the Lefschetz strings and the first two
assertions.

For the exact constant, write \(z=\sum_{|T|=t}u_Tx_T\).
From \(Fz=0\) one obtains, for every \(J\subset I\) with
\(|J|<t\),
\[
 \sum_{\substack{T\supset J\\|T|=t}}u_T=0.
\]
This follows by looking at the coefficient of \(x_J\) in
\(F^{t-|J|}z\), which is \((t-|J|)!\) times the sum. If
\(K\subset I\) has size \(m-t\) and \(R_0=I\setminus K\),
inclusion--exclusion gives
\[
 \sum_{\substack{T\subset K\\|T|=t}}u_T
 =\sum_{J\subset R_0}(-1)^{|J|}
       \sum_{\substack{T\supset J\\|T|=t}}u_T
 =(-1)^tu_{R_0}.
\]
The coefficient of \(x_K\) in \(E^{m-2t}z\) is
\((m-2t)!\) times the left side. This proves
\MathEqref{eq:star-identity}, including \(t=0\).
\end{proof}

\subsection{Proof of the arithmetic Lefschetz reduction}
\label{sec:lefschetz-proof}

We prove \Cref{thm:main-lefschetz}. Suppose first that all the
noncentral cores vanish. On a remaining block put \(k=|S|=2c-1\),
\(m=n-k\), and \(t=r-c\). Then
\[
 n+1-2r=m-2t,
\]
and the desired map is
\(1\otimes E^{m-2t}\). It is an isomorphism by
\Cref{lem:square-zero}. If \(t<0\), both its source and target
in this block vanish. Taking the direct sum proves (ii).

Conversely, suppose \(U_S^c(E)\ne0\) and put
\(d=k+1-2c\). The homologically trivial groups force
\(1\le c\le k\). If \(d>0\), take \(r=c\). The required
map has exponent \(m+d>m\), so it kills
\(U_S^c(E)\otimes1\); it cannot be injective. If \(d<0\),
take \(r=k+1-c\). This satisfies
\(1\le r\le\lfloor(n+1)/2\rfloor\). The target has the
nonzero component \(U_S^c(E)\otimes x_I\) in degree \(c+m\),
whereas the source in this same block has polynomial degree
\(r-c=d<0\) and is zero. Multiplication by \(\ell\) preserves
the block, so surjectivity fails. Hence \(d=0\) for every nonzero
core. This proves (i).

\subsection{Primitive heights and their exact constants}
\label{sec:primitive-proof}

Adjunction and \({}^tP_{S,T,U}=P_{S,U,T}\) show that Chow blocks
are height-orthogonal unless their point and fundamental-class
indices are complementary and their \(S\) indices agree. In the
complementary case,
\begin{equation}\label{eq:complement-height}
 h_X(j_Tu,j_{I\setminus T}v)=h_{Y_S}(u,v).
\end{equation}
Indeed the transpose of the first embedding composed with the
second integrates one degree-one point in each omitted factor.

Define
\[
 P_\ell^r(X_E)=\ker\bigl(
   \ell^{n+2-2r}:V^r(X_E)\to V^{n+2-r}(X_E)\bigr)
\]
for \(1\le r\le\lfloor(n+1)/2\rfloor\).

\begin{theorem}[Primitive sum of squares]
\label[theorem]{thm:primitive-squares}
Let \(|S|=2c-1\), \(m=n-|S|\), and \(0\le t=r-c\le m/2\).
Write a vector in the corresponding block as
\[
 z_S=\sum_{|T|=t}u_Tx_T,
 \qquad u_T\in U_S^c(E).
\]
It is primitive exactly when \(Fz_S=0\). For such a vector,
\begin{equation}\label{eq:primitive-squares}
 q_{\ell,r}(z_S):=(-1)^rh_X(z_S,\ell^{n+1-2r}z_S)
   =(m-2t)!\,a_I\sum_{|T|=t}q_{S,c}(u_T),
\end{equation}
where \(a_I=\prod_{i\in I}a_i\) and
\(q_{S,c}(u)=(-1)^ch_{Y_S}(u,u)\).
\end{theorem}

\begin{proof}
The primitive exponent is \(m-2t+1\), so the equivalence follows
from \Cref{lem:square-zero}. The same lemma gives
\(\ell^{n+1-2r}z_S=(-1)^t(m-2t)!\,*z_S\).
In pairing \(z_S\) with \(*z_S\), only the identical coefficient
indices survive by \MathEqref{eq:complement-height}. Their weights
multiply to \(a_Ta_{I\setminus T}=a_I\). The remaining sign is
\((-1)^{r+t}=(-1)^c\), since \(r=c+t\). This proves the formula.
\end{proof}

\begin{corollary}\label[corollary]{cor:positivity-propagation}
Suppose all the noncentral subproduct cores vanish and all central
forms \(q_{S,c}\) are real positive definite over a fixed number
field. Then the complementary height pairings on \(X\) are
nondegenerate over that field, also after real extension, and
\(q_{\ell,r}\) is real positive definite on \(P_\ell^r\).
Its associated bilinear form \(b_{\ell,r}\) satisfies
\[
 |b_{\ell,r}(z,w)|^2
      \le q_{\ell,r}(z)q_{\ell,r}(w).
\]
If noncentral vanishing holds over every finite extension and each
central integral image has bounded-degree Northcott, then so does
\(\Lambda_X^r(E)\cap P_\ell^r\) for \(q_{\ell,r}\).
\end{corollary}

\begin{proof}
The decomposition into complementary blocks and the definiteness
of their central forms prove nondegeneracy. Distinct \(S\) are
orthogonal. Formula \MathEqref{eq:primitive-squares} therefore
proves primitive definiteness. Cauchy--Schwarz follows by applying
positivity to \(q_{\ell,r}(z+tw)\), \(t\in\R\).

For the bounded-degree assertion, central Northcott gives
nonnegativity over each finite extension: a nonzero rational class
of nonpositive height would have infinitely many integral multiples
in a bounded height ball. For an integral primitive \(z\) of
height at most \(B\), every
coefficient satisfies
\[
 q_{S,c}(u_T)\le
   \frac{B}{(m-2t)!a_I},\qquad
 N_a u_T\in\Lambda_{U_S^c}(E).
\]
Thus central Northcott and \MathEqref{eq:coordinate-denominator}
give finitely many possibilities for every coefficient. There are
only finitely many blocks, and the decomposition is injective.
\end{proof}
